\chapter{ДОМЕННАЯ ПАРАМЕТРИЗАЦИЯ И ЭКСПЕРИМЕНТАЛЬНАЯ СРЕДА}\label{ch:ch6}

В предыдущих главах были сформулированы математическая модель системы автономных мобильных роботов с поддержкой периферийных вычислений, аукционный слой распределения задач, обучаемый контур управления операционными режимами периферийных узлов, а также гибридный метод AMPLE-AMR и его кластерный вариант C-AMPLE-AMR. Цель настоящей главы состоит в задании доменной параметризации, необходимой для последующей реализации программного прототипа и экспериментальной проверки предложенного метода.

Валидационный домен должен проверять не только способность метода уменьшать задержку обработки задач, но и те свойства, ради которых в работе вводится гибридная архитектура: распределение общего дефицитного ресурса, учёт внешнего эффекта назначений, адаптация к долговременной динамике очередей, устойчивость к перегрузке и возможность масштабирования. Поэтому в качестве основного домена выбран роботизированный склад. Такая среда является локальной, допускает размещение периферийных вычислительных узлов вблизи зоны работы роботов, характеризуется неоднородным потоком задач и позволяет естественно моделировать пиковую нагрузку, ухудшение связи, частичную недоступность узлов и рост числа роботов~\cite{Fragapane2021AMRPlanning,Lackner2024AMRIntralogistics,Tahir2025EdgeRoboticsSurvey}.

Роботизированный склад рассматривается как частный, но практически значимый представитель более широкого класса локальных AMR-систем. К этому классу относятся внутрипроизводственная логистика, распределённые участки сборки, больничная логистика, терминальные комплексы и кампусная доставка. Поэтому параметры главы ориентированы на складскую среду, но сама структура модели не зависит от конкретной складской планировки.

В главе решаются следующие задачи:
\begin{enumerate}
    \item определить границы складской модели и принципы доменной параметризации;
    \item задать параметры автономных мобильных роботов и классы вычислительных задач;
    \item определить классы периферийных вычислительных узлов и их операционные режимы;
    \item задать модель сетевой связности, задержек и ухудшения связи;
    \item определить профили нагрузки и сценарии проверки;
    \item задать размеры экспериментальных конфигураций и критерии реалистичности среды.
\end{enumerate}

Экспериментальные результаты в настоящей главе не приводятся. Они формируются программным прототипом, описанным в главе~\ref{ch:ch7}, и анализируются в главе~\ref{ch:ch8}.

\section{Принципы доменной параметризации}

Доменная параметризация строится на трёх принципах.

Первый принцип состоит в отделении проверяемых свойств метода от случайных деталей реализации склада. Метод AMPLE-AMR управляет распределением вычислительных задач и режимами периферийных узлов, но не решает задачу построения траекторий роботов. Поэтому в модели сохраняются те элементы складской среды, которые влияют на вычислительную разгрузку: число роботов, поток задач, достижимость узлов, сетевые задержки, очереди, вычислительные ресурсы и отказы. Детальная геометрия стеллажей и локальное планирование движения не включаются в базовую модель.

Второй принцип состоит в проверяемой связи параметров с прикладным доменом. Значения скорости, полезной нагрузки и времени работы роботов используются как ориентиры для выбора масштаба складской системы, а не как непосредственные переменные целевой функции. Параметры периферийных узлов задаются через нормированные вычислительные ёмкости, соответствующие слабым, средним и мощным узлам, поскольку задача работы состоит не в сравнении конкретных микропроцессорных архитектур, а в проверке метода распределения при гетерогенных ресурсах.

Третий принцип состоит в воспроизводимости. Все параметры, влияющие на эксперимент, должны задаваться явно: размеры конфигурации, интенсивность поступления задач, распределение классов задач, параметры задержек, вероятности отказов, длительность эпизода и начальные состояния очередей. Это позволяет повторять эксперименты и отделять влияние метода распределения от влияния случайной генерации сценариев.

\section{Роботизированный склад как валидационный домен}

Роботизированный склад представляет собой локальную среду, в которой группа автономных мобильных роботов перемещает грузы между зонами хранения, комплектации, упаковки, приёмки и отгрузки. Роботы взаимодействуют с системой управления складом, точками беспроводного доступа, зонами зарядки и локальной вычислительной инфраструктурой. В такой среде вычислительные задачи возникают не изолированно, а как часть непрерывного операционного процесса: робот должен воспринимать окружение, уточнять своё положение, обмениваться телеметрией, обновлять маршрут и выполнять вспомогательные задачи контроля.

С точки зрения настоящей работы складской домен обладает следующими свойствами.

Во-первых, склад является локальной ограниченной областью. Это позволяет разместить периферийные узлы вблизи рабочих зон и использовать их как общий вычислительный ресурс для группы роботов.

Во-вторых, поток задач неоднороден. Одни задачи связаны с безопасностью движения и имеют жёсткие ограничения по задержке, другие относятся к инвентаризации, диагностике или обновлению карты и допускают более мягкие сроки выполнения.

В-третьих, нагрузка изменяется во времени. Пиковые волны заказов, периоды комплектации, инвентаризация, смена рабочей зоны и частичное отключение оборудования создают нестационарный поток задач.

В-четвёртых, назначение одной задачи влияет на обслуживание остальных задач. Если ресурсоёмкая задача занимает периферийный узел, то растёт очередь и ухудшается доступность этого узла для следующих задач. Это является прямым проявлением внешнего эффекта, для оценки которого в работе используется аукционный слой.

В-пятых, с ростом числа роботов и периферийных узлов глобальная задача выбора распределения становится вычислительно дорогой. Поэтому складской домен позволяет проверять не только базовый метод AMPLE-AMR, но и кластерный вариант C-AMPLE-AMR.

Обобщённая схема складского сценария приведена на рис.~\ref{fig:ch6_warehouse_scenario}.

\begin{figure}[h]
    \centering
    \includegraphics[width=0.95\textwidth]{ch6_warehouse_scenario}
    \caption{Складская среда с автономными мобильными роботами, периферийными узлами и зонами изменения нагрузки}
    \label{fig:ch6_warehouse_scenario}
\end{figure}

\section{Модель складской среды}

Складская среда моделируется как прямоугольная область
\[
    \mathcal{W} = [0,L_x] \times [0,L_y],
\]
где расположены рабочие зоны, зоны хранения, зоны комплектации, точки доступа, периферийные вычислительные узлы и автономные мобильные роботы. В базовой конфигурации размеры области выбираются из диапазона
\[
    L_x \in [60,200] \text{ м}, \qquad L_y \in [40,120] \text{ м}.
\]
Эти значения задают масштаб локальной складской среды и используются для моделирования сетевых задержек и достижимости узлов, но не требуют детального моделирования каждой полки или прохода.

Связность системы задаётся графом
\[
    G(t)=(\mathcal{V},\mathcal{E}(t),w(t)),
\]
где вершинами являются роботы, точки доступа и периферийные узлы, а рёбра отражают достижимость, задержку передачи, пропускную способность и деградацию связи. Такая графовая модель соответствует уровню абстракции задачи: AMPLE-AMR распределяет вычислительные задачи и управляет режимами узлов, а не строит траектории движения роботов.

Рабочая область делится на зоны:
\begin{enumerate}
    \item зоны стабильного движения, где сетевые параметры близки к номинальным;
    \item зоны интенсивного движения, где возрастает вероятность пикового поступления задач;
    \item зоны ухудшенной связи, где задержки и потери пакетов выше из-за препятствий, металлических конструкций или перегрузки беспроводной сети;
    \item зоны обслуживания и зарядки, где часть роботов может временно не формировать вычислительные задачи;
    \item зоны отказа или обслуживания инфраструктуры, в которых отдельные периферийные узлы становятся недоступными.
\end{enumerate}

Такое разбиение позволяет формировать сценарии, необходимые для проверки основных свойств метода: устойчивости к пикам нагрузки, ухудшению связи, отказам узлов и росту числа роботов.

\section{Параметры автономных мобильных роботов}

Множество роботов обозначается
\[
    \mathcal{R}=\{r_1,\ldots,r_N\}.
\]
Каждый робот характеризуется положением $\mathbf{p}_i(t)$, скоростью движения, текущим состоянием связи, уровнем занятости и потоком вычислительных задач. В рамках настоящей работы физические параметры робота используются для задания реалистичного масштаба домена, а не как основной объект оптимизации.

В качестве ориентиров используются промышленные автономные мобильные роботы, применяемые во внутрискладской и производственной логистике. Например, MiR250 имеет полезную нагрузку до 250~кг, максимальную скорость 2{,}0~м/с и время активной работы до 17~ч~30~мин при использовании батареи увеличенной ёмкости~\cite{MiR250Spec2026}. OTTO~100 рассчитан на перемещение грузов до 150~кг; в технических спецификациях также указывается максимальная скорость 2{,}0~м/с и поддержка беспроводной связи Wi-Fi 2{,}4/5~ГГц~\cite{OTTO100Spec2026} (см. табл.~\ref{tab:amr_real_params} в приложении А).

В симуляционной среде используются следующие параметры роботов:
\begin{enumerate}
    \item число роботов $N$;
    \item положение робота $\mathbf{p}_i(t)$;
    \item класс рабочей зоны;
    \item состояние связи с точками доступа;
    \item интенсивность генерации задач $\lambda_i(t)$;
    \item распределение классов задач;
    \item признак локальной перегрузки или временного обслуживания.
\end{enumerate}

Параметры полезной нагрузки и батареи не входят в базовую функцию общественного благосостояния. Это сделано для сохранения предметной направленности работы: исследуется распределение вычислительных задач, а не совместная задача маршрутизации, зарядки и распределения груза. При расширении модели энергетическое состояние робота может быть добавлено в функцию полезности или в ограничения допустимости разгрузки.

\section{Классы вычислительных задач}

Каждый робот формирует поток задач
\[
    \mathcal{T}(t)=\{\tau_1,\ldots,\tau_{M(t)}\}.
\]
Тип задачи задаётся вектором
\[
    \theta_{\tau}=\bigl(c_{\tau}, d_{\tau}, \rho_{\tau}, q^{cpu}_{\tau}, q^{mem}_{\tau}, q^{gpu}_{\tau}, b^{in}_{\tau}, b^{out}_{\tau}\bigr),
\]
где $c_{\tau}$ --- класс задачи, $d_{\tau}$ --- допустимый срок завершения, $\rho_{\tau}$ --- приоритет, $q^{cpu}_{\tau}$, $q^{mem}_{\tau}$ и $q^{gpu}_{\tau}$ --- требования к вычислительным ресурсам, $b^{in}_{\tau}$ и $b^{out}_{\tau}$ --- объёмы входных и выходных данных.

В базовой конфигурации используются пять классов задач:
\begin{enumerate}
    \item задачи безопасности и быстрой реакции;
    \item задачи локализации и уточнения положения;
    \item задачи восприятия и распознавания объектов;
    \item задачи обновления карты и координации движения;
    \item задачи инвентаризации, диагностики и пакетной обработки журналов.
\end{enumerate}

Для задач безопасности допускается бортовой резерв исполнения, поскольку потеря связи с периферийной инфраструктурой не должна приводить к потере безопасного поведения робота. Для задач восприятия, построения карты и инвентаризации допускается разгрузка на периферийные узлы. Для диагностических задач допускается более мягкий срок выполнения, но они могут создавать значительную нагрузку при массовом поступлении.

Нормированные параметры классов задач приведены в табл.~\ref{tab:ch6_task_classes}. Таблица задаёт не физические единицы конкретного устройства, а относительные значения, используемые симулятором для расчёта времени обработки, сетевой стоимости и функции полезности.

\begin{table}[h]
\centering
\caption{Нормированные параметры классов вычислительных задач}
\label{tab:ch6_task_classes}
\begin{tabularx}{\textwidth}{X c c c c}
\toprule
\textbf{Класс задачи} & \textbf{Срок, мс} & \textbf{CPU} & \textbf{GPU} & \textbf{Приоритет} \\
\midrule
Безопасность и быстрая реакция & 50--150 & 1--2 & 0--1 & высокий \\
Локализация и уточнение положения & 100--300 & 1--3 & 0--2 & высокий \\
Восприятие и распознавание объектов & 150--500 & 2--5 & 2--6 & средний/высокий \\
Обновление карты и координация & 300--1000 & 2--6 & 1--4 & средний \\
Инвентаризация и диагностика & 1000--5000 & 1--8 & 0--6 & низкий/средний \\
\bottomrule
\end{tabularx}
\end{table}

Чем меньше допустимый срок и выше приоритет, тем больше предельная полезность своевременного выполнения задачи. Если задача завершена после срока $d_{\tau}$, её вклад в общественное благосостояние уменьшается штрафной функцией, определённой во второй главе. Это позволяет проверять не только среднюю задержку, но и способность метода защищать критические задачи от вытеснения менее срочными задачами.

\section{Периферийная вычислительная инфраструктура}

Множество периферийных узлов обозначается
\[
    \mathcal{E}=\{e_1,\ldots,e_K\}.
\]
Каждый узел характеризуется вычислительной ёмкостью, объёмом памяти, наличием ускорителя, очередью задач, состоянием доступности, положением в складской среде и текущим операционным режимом.

В качестве аппаратных ориентиров используются семейства встраиваемых вычислительных платформ для робототехники и периферийного искусственного интеллекта. Согласно спецификациям NVIDIA, серия Jetson Orin Nano обеспечивает до 67~трлн операций в секунду для задач искусственного интеллекта, серия Jetson Orin NX --- до 157~трлн операций в секунду, а Jetson AGX Orin --- до 275~трлн операций в секунду при настраиваемом энергопотреблении~\cite{NvidiaJetsonOrin2026}. В симуляционной среде эти значения не используются напрямую, но задают реалистичное соотношение слабых, средних и мощных периферийных узлов.

В базовой модели используются три класса узлов:
\begin{enumerate}
    \item \textbf{Edge-S} --- малый периферийный узел, размещаемый вблизи точки доступа или локальной рабочей зоны;
    \item \textbf{Edge-M} --- средний периферийный узел, обслуживающий несколько рабочих зон;
    \item \textbf{Edge-G} --- мощный узел с графическим ускорителем, используемый для задач восприятия, построения карты и пакетной обработки.
\end{enumerate}

Их нормированные ёмкости задаются как
\[
    F_{S}=(1,1,1), \qquad F_{M}=(2,2,3), \qquad F_{G}=(4,4,8),
\]
где компоненты соответствуют относительным ресурсам центрального процессора, памяти и ускорителя. Физическая ёмкость узла не равна ёмкости, доступной для новых задач. В текущем раунде аукционный слой использует экспонируемую ёмкость
\[
    \widetilde{F}_{j}(t)=\alpha_j(a_j(t),s_j(t))F_j,
\]
где $a_j(t)$ --- выбранный операционный режим, $s_j(t)$ --- состояние узла, а $\alpha_j(\cdot)$ --- коэффициент открытия ресурсов для новых задач.

Такое различие принципиально для AMPLE-AMR. Обучаемый контур не меняет ценность задач и не подменяет аукционный механизм, а изменяет условия распределения через экспонируемую ёмкость, резервирование и правила допуска задач.

\section{Операционные режимы периферийных узлов}

Операционный режим периферийного узла определяет, какую часть ресурсов узел готов открыть для новых задач и какие ограничения применяются к очереди обслуживания. В базовой конфигурации задаётся конечное множество режимов
\[
    \mathcal{A}_j=\{a^{normal},a^{reserve},a^{critical},a^{protect},a^{recovery}\}.
\]

Режим $a^{normal}$ соответствует штатной работе: узел принимает задачи всех допустимых классов в пределах текущей экспонируемой ёмкости. Режим $a^{reserve}$ резервирует часть ресурсов под будущие критические задачи. Режим $a^{critical}$ допускает преимущественно задачи высокого приоритета. Режим $a^{protect}$ используется при перегрузке: узел уменьшает экспонируемую ёмкость и ограничивает приём менее срочных задач. Режим $a^{recovery}$ соответствует восстановлению после отказа или деградации, когда узел временно принимает только ограниченный поток задач.

Выбор именно этих режимов связан с целями проверки метода. Если режимы отсутствуют, повторяемый аукционный слой каждый раз распределяет задачи в рамках текущей видимой ёмкости, но не управляет резервированием и защитой от будущей перегрузки. Если режимы есть, обучаемая политика может заранее ограничить приём менее срочных задач, сохранить ресурс для критического потока или снизить нагрузку на узел с ухудшающейся связью.

\section{Модель сетевой связности и задержек}

Для каждой пары робот--периферийный узел задаётся сетевой путь через точку доступа. Сетевая задержка разгрузки задачи $\tau$ с робота $r_i$ на узел $e_j$ записывается как
\[
    T^{net}_{\tau ij}(t)=T^{up}_{\tau ij}(t)+T^{queue,net}_{ij}(t)+T^{down}_{\tau ij}(t),
\]
где $T^{up}_{\tau ij}(t)$ --- время передачи входных данных, $T^{queue,net}_{ij}(t)$ --- задержка из-за сетевой очереди и занятости канала, $T^{down}_{\tau ij}(t)$ --- время передачи результата.

В базовой конфигурации задержка канала представляется как
\[
    T^{link}_{ij}(t)=T^{base}_{ij}\cdot \gamma_{ij}(t)+\xi_{ij}(t),
\]
где $T^{base}_{ij}$ --- базовая задержка для зоны, $\gamma_{ij}(t)$ --- коэффициент деградации связи, а $\xi_{ij}(t)$ --- случайная добавка. Коэффициент деградации увеличивается в зонах ухудшенной связи, при перегрузке точки доступа и при частичной потере достижимости.

Пара робот--узел считается недоступной, если нарушено хотя бы одно условие:
\[
    T^{link}_{ij}(t) > T^{max}_{link}, \qquad B_{ij}(t)<B^{min}, \qquad z_{ij}(t)=0,
\]
где $B_{ij}(t)$ --- доступная пропускная способность, $z_{ij}(t)$ --- бинарный признак достижимости. Такой подход согласуется с работами по вычислительной разгрузке мобильных роботов, где соединения исключаются из множества допустимых при нарушении порогов задержки или пропускной способности~\cite{Baruffa2025AMRResourceAssignment,Djordjevic2023GroundRoboticsOffloading}.

В стресс-сценариях сетевые параметры изменяются во времени. Это позволяет проверить, способен ли метод не только выбирать ближайший или самый мощный узел, но и адаптироваться к изменению доступности инфраструктуры (см. табл.~\ref{tab:network_params} в приложении А).

\section{Профили нагрузки}

Поток задач формируется как нестационарный процесс. Для робота $r_i$ интенсивность поступления задач обозначается $\lambda_i(t)$. В базовой модели используются следующие профили нагрузки:
\begin{enumerate}
    \item \textbf{стабильный профиль}: интенсивность задач близка к постоянной, отказы и ухудшение связи отсутствуют;
    \item \textbf{пиковый профиль}: на ограниченном интервале времени возрастает интенсивность задач восприятия, локализации и инвентаризации;
    \item \textbf{смешанный профиль}: в разных зонах склада одновременно возникают задачи разных классов с различными сроками и приоритетами;
    \item \textbf{профиль ухудшения связи}: часть роботов временно теряет доступ к ближайшим периферийным узлам или получает увеличенную задержку;
    \item \textbf{профиль отказов}: один или несколько узлов временно становятся недоступными или переходят в режим восстановления;
    \item \textbf{профиль масштабирования}: число роботов, задач и узлов постепенно увеличивается для проверки вычислительных накладных расходов.
\end{enumerate}

Профили нагрузки подбираются так, чтобы проверять разные элементы метода. Стабильный профиль нужен для оценки базовой эффективности распределения. Пиковый профиль проверяет наличие внешнего эффекта и способность метода защищать критические задачи. Ухудшение связи и отказы проверяют адаптацию режимов узлов. Профиль масштабирования проверяет необходимость перехода от глобального аукциона к кластерной версии.

\section{Размеры экспериментальных конфигураций}

Для проверки масштабируемости используются три группы конфигураций:
\begin{enumerate}
    \item малая конфигурация: $N=10$--$20$ роботов, $K=3$--$5$ периферийных узлов;
    \item средняя конфигурация: $N=30$--$60$ роботов, $K=6$--$10$ периферийных узлов;
    \item крупная конфигурация: $N=80$--$150$ роботов, $K=12$--$20$ периферийных узлов.
\end{enumerate}

Малая конфигурация используется для отладки модели и проверки свойств аукционного слоя. Средняя конфигурация является основной для сравнения методов. Крупная конфигурация используется для проверки вычислительных накладных расходов, устойчивости к росту числа задач и целесообразности применения C-AMPLE-AMR.

Длительность одного эпизода задаётся числом управляющих шагов
\[
    H \in \{500,1000,2000\}.
\]
На каждом шаге формируется множество новых задач, обновляются очереди, сетевые параметры и состояния узлов. После этого обучаемая политика выбирает операционные режимы узлов, а аукционный слой распределяет допустимые задачи в рамках экспонируемой ёмкости.

\section{Сценарии проверки метода}

В главе~\ref{ch:ch8} экспериментальная проверка должна быть построена не как единичное сравнение средних задержек, а как набор сценариев, каждый из которых проверяет отдельное утверждение диссертации.

Первый сценарий проверяет пользу аукционного слоя. В нём сравниваются эвристическое распределение и распределение на основе механизма Викри--Кларка--Гровса при фиксированных операционных режимах узлов. Проверяемое утверждение: учёт общественного благосостояния и внешнего эффекта должен улучшать обслуживание неоднородного потока задач по сравнению с простыми правилами назначения.

Второй сценарий проверяет пользу адаптации. В нём сравниваются фиксированные режимы узлов и обучаемое управление режимами при одинаковом аукционном слое. Проверяемое утверждение: управление экспонируемой ёмкостью и резервированием ресурсов должно уменьшать перегрузку и долю просроченных критических задач в динамической среде.

Третий сценарий проверяет пользу гибридизации. В нём сравниваются методы, использующие только аукцион, только обучаемую эвристику и полный метод AMPLE-AMR. Проверяемое утверждение: разделение локального распределения и долгосрочной адаптации должно обеспечивать более устойчивое качество обслуживания, чем использование каждого компонента отдельно.

Четвёртый сценарий проверяет кластерную версию. В нём глобальный аукцион сравнивается с C-AMPLE-AMR при росте числа роботов и узлов. Проверяемое утверждение: локализация распределения внутри кластеров уменьшает вычислительные и коммуникационные накладные расходы при контролируемой потере качества распределения.

Пятый сценарий проверяет диагностическую ценность аукционных сигналов. В нём анализируется связь VCG-платежей, оценок внешнего эффекта, загрузки узлов, длины очередей и частоты отклонения задач. Проверяемое утверждение: внутренние теневые цены должны возрастать в состояниях дефицита ресурса и тем самым служить информативными признаками для обучаемого контура.

\section{Критерии реалистичности и ограничения модели}

Реалистичность экспериментальной среды обеспечивается не точным воспроизведением конкретного склада, а сохранением тех факторов, которые принципиальны для задачи распределения вычислительных задач: гетерогенность роботов и задач, ограниченность периферийных узлов, изменчивая связь, очереди обслуживания, нестационарная нагрузка, отказы и масштабирование.

Модель имеет следующие ограничения. Во-первых, траектории роботов не оптимизируются совместно с вычислительной разгрузкой. Во-вторых, энергетическое состояние роботов не входит в базовый критерий общественного благосостояния. В-третьих, сетевые параметры задаются абстрактной графовой моделью, а не детальной физической моделью распространения радиосигнала. В-четвёртых, аппаратные характеристики периферийных узлов нормируются и не претендуют на точное моделирование конкретных процессоров.

Эти ограничения являются осознанными. Они позволяют изолировать предмет исследования: адаптивное распределение вычислительных задач автономных мобильных роботов между периферийными узлами. При необходимости модель может быть расширена за счёт включения маршрутизации, зарядки, детальной радиомодели и энергетического состояния роботов.

\section{Выводы по главе}

В главе задана доменная параметризация экспериментальной среды для проверки метода AMPLE-AMR и его кластерного варианта C-AMPLE-AMR. В качестве основного валидационного домена выбран роботизированный склад, поскольку он сочетает локальность инфраструктуры, общий дефицитный вычислительный ресурс, неоднородный поток задач, динамическую нагрузку, изменчивую связь и возможность масштабирования.

Параметры роботов, периферийных узлов, вычислительных задач, сетевой связности, операционных режимов и профилей нагрузки заданы так, чтобы проверять основные утверждения диссертации: пользу аукционного распределения, необходимость оценки внешнего эффекта, пользу адаптивного управления режимами узлов, устойчивость к деградации инфраструктуры и целесообразность кластерной версии при росте размера системы.

Полученная параметризация используется в главе~\ref{ch:ch7} при описании программного прототипа и в главе~\ref{ch:ch8} при формировании программы экспериментальной валидации.
